\documentclass[UTF8,a4paper,11pt,fontset=none]{ctexart}
\usepackage[top=5.3mm,bottom=14.4mm,left=9.8mm,right=9.8mm]{geometry}
\usepackage{fontspec}
\usepackage{xeCJK}
\usepackage{enumitem}
\usepackage{tabularx}
\usepackage{titlesec}
\usepackage[hidelinks]{hyperref}
\usepackage{xcolor}
\usepackage{microtype}
\usepackage{setspace}
\usepackage{ragged2e}

\setmainfont{Noto Serif}[BoldFont=Noto Serif SemiBold]
\setCJKmainfont[BoldFont=Noto Serif CJK SC Bold]{Noto Serif CJK SC}
\setsansfont{Noto Sans}[BoldFont=Noto Sans Bold]
\setCJKsansfont[BoldFont=Noto Sans CJK SC Bold]{Noto Sans CJK SC}
\setmonofont{DejaVu Sans Mono}
\setCJKmonofont{Noto Sans Mono CJK SC}

\pagestyle{empty}
\setlength{\parindent}{0pt}
\setlength{\tabcolsep}{0pt}
\setlength{\parskip}{0pt}
\urlstyle{same}
\raggedbottom

\titleformat{\section}{\sffamily\bfseries\fontsize{12.2pt}{13.8pt}\selectfont}{}{0pt}{}[\vspace{0.3pt}{\color{black!55}\titlerule[0.35pt]}]
\titlespacing*{\section}{0pt}{0.55pt}{0.05pt}

\newcommand{\roleline}[2]{%
  \begin{tabularx}{\textwidth}{@{}Xr@{}}
    {\sffamily\bfseries\fontsize{10.78pt}{12.05pt}\selectfont #1} & {\sffamily\fontsize{8.8pt}{10.2pt}\selectfont #2}
  \end{tabularx}\vspace{-0.05pt}
}

\newcommand{\subroleline}[2]{%
  \vspace{-0.35pt}%
  \begin{tabularx}{\textwidth}{@{}Xr@{}}
    {\sffamily\bfseries\fontsize{10.2pt}{11.45pt}\selectfont #1} & {\sffamily\fontsize{8.65pt}{10.1pt}\selectfont #2}
  \end{tabularx}\vspace{-0.2pt}
}

\newcommand{\resumeliststart}{\begin{itemize}[leftmargin=1.42em,label=\textbullet,itemsep=0pt,parsep=0pt,topsep=-0.05pt,partopsep=0pt]}
\newcommand{\resumeitem}[1]{\item\fontsize{10.12pt}{10.68pt}\selectfont\RaggedRight #1}
\newcommand{\resumelistend}{\end{itemize}\vspace{0pt}}

\begin{document}

\begin{center}
  {\sffamily\bfseries\fontsize{18.6pt}{20.5pt}\selectfont Leo Li}\\[0.65pt]
  {\sffamily\fontsize{8.35pt}{9.35pt}\selectfont
  +1 778-779-2889 \;|\;
  \href{mailto:cheerleaderleo@outlook.com}{cheerleaderleo@outlook.com} \;|\;
  \href{https://www.linkedin.com/in/leoooli/}{linkedin.com/in/leoooli} \;|\;
  \href{https://github.com/teamleaderleo}{github.com/teamleaderleo} \;|\;
  \href{https://teamleaderleo.com/}{teamleaderleo.com}}\\[0.55pt]
  {\sffamily\bfseries\fontsize{9.55pt}{11.2pt}\selectfont 求职方向：性能优化 / 游戏运行时与客户端基础技术 / 跨平台工具链}
\end{center}
\vspace{-4.0pt}

\section{核心项目}
\roleline{Preflight — 游戏运行时性能优化与跨平台启动工具（个人开源项目）}{2026.07 -- 至今}
{\sffamily\fontsize{9.0pt}{10.08pt}\selectfont\href{https://github.com/teamleaderleo/preflight}{github.com/teamleaderleo/preflight}}\vspace{0.15pt}
\resumeliststart
  \resumeitem{\textsf{\bfseries 整体性能：}分析经过混淆的 JVM 游戏运行时及 \textbf{83 个第三方 Mod} 的启动、资源加载与运行时链路，结合 JFR、运行时插桩、结果缓存、预计算与字节码改写，将启动时间从 \textbf{112.17s 降至 13.69s（8.19×）}。}
  \resumeitem{\textsf{\bfseries 数据加载：}将 5 套加载器中重复的 JSON/CSV 读取与合并收敛到共享解析与读取缓存层，覆盖 \textbf{39,017 次 JSON 调用 / 8,378 条路径}，使 \texttt{SpecStore} 从 \textbf{19.8s 降至 9.8s}、合并读取耗时从 \textbf{2.172s 降至 0.300s}。}
  \resumeitem{\textsf{\bfseries 资源与存储：}将纹理缓存命中判断移到曾阻塞启动约 \textbf{27s} 的单线程预取队列之前，并移除历史纹理上传中的约 \textbf{1.22 GiB 无效显存填充}；重做可重建中间产物的持久化与单包写入路径，将准备时间从 \textbf{200.77s 降至 16.21s}、稳态存储从 \textbf{4.76 GB 降至约 1.1 GB}。}
  \resumeitem{\textsf{\bfseries 存储与 I/O：}在逻辑纹理内容不变的情况下，按真实启动访问顺序重排物理存储，使同一纹理集的启动时间从 \textbf{33.53s 降至 14.174s}；另以基于变更跟踪的增量索引替代高频全局 O(n) 扫描，消除 \textbf{7,910 万次实体引用检查}并短路 \textbf{1.179 亿次}未变化的重复计算。}
  \resumeitem{\textsf{\bfseries 运行时：}缓存 \textbf{228 次 Janino 动态编译请求}，再将 \textbf{36,332 个生成类条目收敛为 280 个唯一类}，使生成类映射从 \textbf{145.96 MiB 降至 1.13 MiB}、重放从 \textbf{1.501s 降至 29ms}。}
  \resumeitem{\textsf{\bfseries 产品化：}将 Java 性能核心做成 Windows / macOS / Linux 桌面应用，采用 React UI + Rust/Tauri host，内置 Java 运行时，并实现启动/游玩记录、诊断、签名更新与回滚。}
\resumelistend

\section{开源工程}
\subroleline{Cloud Hypervisor — Rust / VMM / 虚拟化与存储}{2026.08 · 4 merged}
\resumeliststart
  \resumeitem{\textsf{\bfseries 生命周期与故障处理：}修复 VM 关闭竞态，改为等待 VMM 的真实 \texttt{shutdown} 事件后再复用 VM/磁盘；将 ACPI 表构建中的地址溢出、缺失 \texttt{fw\_cfg}、guest-memory write 等失败转换为可传播的 VM 启动错误，避免 VMM panic。}
  \resumeitem{\textsf{\bfseries 设备与存储正确性：}修复 VFIO sparse BAR 映射边界检查，拒绝跨越未映射区间的 DMA 请求；修复 QCOW L2 元数据所有权与失败顺序，避免仍被镜像引用的元数据在 reopen 后被误判为空闲空间并再次分配。}
\resumelistend

\subroleline{Vercel AI SDK — TypeScript / Web Runtime}{2026.08 · 1 merged + 2 项方案采用}
\resumeliststart
  \resumeitem{\textsf{\bfseries 运行时正确性：}修复 global / sticky 正则导致相同 URL 重复检查返回不同结果的问题；修复 Web Stream 读取失败后的 reader 清理，使失败读取不会遗留 locked stream；并确保下载超限错误不会被后续 \texttt{reader.cancel()} 清理失败覆盖。}
\resumelistend

\subroleline{Vite / Workers SDK — 构建工具链 / 运行时生命周期}{2026.08 · 4 merged + 1 open}
\resumeliststart
  \resumeitem{\textsf{\bfseries 构建与清理：}修复 Vite 构建失败后跳过插件清理、依赖分析中的临时 Rolldown build 未释放，以及重复配置解析导致已预热依赖缓存被误重建的问题（最后一项 open）。}
  \resumeitem{\textsf{\bfseries 进程与状态：}修复 Miniflare 辅助清理卡住或失败时 \texttt{workerd} 退出被延迟的问题，并避免 Cloudflare Access 凭据被删除或变得不完整后继续使用陈旧认证信息。}
\resumelistend

\section{实习经历}
\subroleline{IBM — Software Developer Intern}{Toronto, Canada · 2021.05 -- 2022.08}
\resumeliststart
  \resumeitem{\textsf{\bfseries 故障处理：}参与 IBM Cloud AI/ML 与数据工作流的 Java 端到端测试与重构，覆盖 Kafka、Spark、Snowflake、混合云与本地部署环境；发现关键 \textbf{RBAC 权限缺陷}并推动三个团队协同完成紧急修复。}
  \resumeitem{\textsf{\bfseries 开发效率：}整理并替换过时的 SDK / runtime 配置流程，将开发环境搭建时间从约 \textbf{3 小时缩短至 15 分钟}。}
\resumelistend

\section{教育背景}
\subroleline{多伦多大学（University of Toronto）｜理学学士｜数学、统计学与计算机科学}{2024.06}

\vspace{-4.2pt}\section{技能}
{\fontsize{9.72pt}{10.64pt}\selectfont
\textbf{编程语言：}Rust、Java、C、Python、TypeScript / JavaScript、Go、SQL\\[0pt]
\textbf{运行时与性能：}Linux、JVM / JFR、Java 字节码插桩与改写、性能分析、缓存与存储布局、生命周期与故障调试\\[0pt]
\textbf{平台与工具：}Windows / macOS / Linux 跨平台开发、Tauri、React / Vite、Docker、Git、CI/CD}

\end{document}
